Seja \(S_n\) o conjunto de permutações da sequência \((1,2,\dots, n)\). Para cada permutação \(\pi=(\pi_1, \dots, \pi_n)\in S_n\), seja \(\mathrm{inv}(\pi)\) o número de pares \(1\le i < j \le n\) com \(\pi_i>\pi_j\); ou seja, o número de inversões em \(\pi\). Denote por \(f(n)\) o número de permutações \(\pi\in S_n\) para as quais \(\mathrm{inv}(\pi)\) é divisível por \(n+1\).

Prove que existem infinitos primos \(p\) tais que \(f(p-1)>\frac{(p-1)!}{p}\), e infinitos primos \(p\) tais que \(f(p-1)<\frac{(p-1)!}{p}\).
